% Elaborado por Fabian Gomez
% Version 1.0

\section{Traceability and Change Management}\label{sec:traceability}

Esta sección define cómo se mantiene la trazabilidad de los requisitos del SyRS/SRS del Rover Olympus (RF/RNF/CON) a través de los artefactos del proyecto (ConOps, Casos de Uso, arquitectura MBSE/SysML, PBS/ASM, N2 y V\&V), y cómo se gestionan cambios para preservar un baseline coherente del SRS final. \cite{omg_sysml_mbse, nasa_cm, segoldmine_ppi}  

Dado que los documentos main/main1/main3 son iteraciones previas del mismo borrador, el objetivo de esta sección es establecer una disciplina de ``baseline único'': (i) Una definición maestra por requisito en la Sección 11. (ii) Referencias cruzadas por ID en ConOps/Interfaces/V\&V/Arquitectura. (iii) Un proceso de control de cambios proporcional al proyecto académico que evite divergencias entre copias del documento y del modelo. \cite{nasa_cm, segoldmine_ppi}

La gestión de configuración/cambios es una práctica de ingeniería de sistemas aplicada a lo largo del ciclo de vida para controlar cambios a características funcionales y físicas del producto, asegurar consistencia entre el producto y su documentación/modelos, y mantener visibilidad del estado real de la configuración.

\textit{Justificación:} Se añade este contenido siguiendo buenas prácticas NASA SE/CM porque los borradores describen trazabilidad, pero no formalizan el proceso mínimo de control de cambios y baseline requerido para mantener integridad del SRS final y su relación con modelos MBSE.

\subsection{Requirements Traceability Strategy}\label{subsec:req-traceability}

Estrategia de trazabilidad (digital thread académico del proyecto): La trazabilidad se mantiene como un hilo continuo desde la intención operacional hasta la evidencia de cumplimiento, con la siguiente cadena principal:
\begin{enumerate}
  \item ConOps y Casos de Uso (UC-01..UC-05) $\rightarrow$ Requisitos (RF/RNF/CON) \cite{segoldmine_ppi, nasa_cm}
  \item Requisitos (RF/RNF/CON) $\rightarrow$ Arquitectura lógica (BDD/IBD de 7 subsistemas)
  \item Arquitectura lógica $\rightarrow$ Arquitectura física (PBS/ASM) y conectividad (N2)
  \item Requisitos $\rightarrow$ Plan V\&V (métodos T/A/D/I) $\rightarrow$ Evidencia (resultados)
\end{enumerate}

\textbf{Regla operativa de baseline:} La Sección 11 contiene la definición maestra de RF/RNF/CON; ConOps, Interfaces, Arquitectura y V\&V solo referencian por ID para evitar múltiples definiciones inconsistentes entre secciones. \cite{nasa_cm, segoldmine_ppi}

\textbf{Enfoque MBSE para trazabilidad:} Se recomienda mantener enlaces trazables entre objetos de requisitos y elementos del modelo SysML (bloques, actividades, estados, interfaces), de modo que el modelo funcione como ``fuente autoritativa'' y habilite revisiones cruzadas (requisitos $\leftrightarrow$ arquitectura $\leftrightarrow$ V\&V).

\textit{Justificación:} Se añade este contenido siguiendo buenas prácticas MBSE/NASA porque los borradores incluyen modelos (BDD/IBD/N2/PBS) pero no explicitaban la política de enlaces trazables como parte de la estrategia formal de trazabilidad.

\subsection{COTS Traceability}\label{subsec:cots-traceability}

\textbf{Objetivo:} Asegurar trazabilidad explícita de la implementación física COTS hacia los requisitos del sistema, de modo que cada función/requisito crítico del UC-01 pueda rastrearse hasta: (i) el ensamble/componente (ASM/PBS), (ii) la interfaz interna (N2), y (iii) la evidencia V\&V correspondiente. \cite{nasa_cm, segoldmine_ppi}

\textbf{Estrategia de trazabilidad COTS (baseline del proyecto):}
\begin{enumerate}
  \item \textbf{Mapeo Función/Requisito $\rightarrow$ Ensamble (ASM-xxxx) $\rightarrow$ Ítem COTS:} El PBS/ASM organiza el hardware en ensamblajes y permite mapear qué componentes implementan cada capacidad. \cite{segoldmine_ppi, nasa_cm}
  \item \textbf{Mapeo Ensamble/Ítem $\rightarrow$ Interfaz interna (N2) $\rightarrow$ Riesgos de integración:} La matriz N2 detalla conexiones (CSI-2, I2C, PWM, UART, DC) y se usa para identificar dependencias e impactos: por ejemplo, RF-006 depende de COMM$\leftrightarrow$C\&DH y de la ruta hasta el enlace; RF-005 depende de GNC/PROP/EPS$\leftrightarrow$C\&DH (y de la cadena GNC$\rightarrow$PROP) para ejecutar el stop en tiempo. \cite{segoldmine_ppi, nasa_cm, omg_sysml_mbse}
  \item \textbf{Mapeo Ítem/Interfaz $\rightarrow$ Método V\&V $\rightarrow$ Evidencia:} La trazabilidad COTS concluye en evidencia verificable: pruebas de adquisición CSI (RF-001), pruebas de clasificación (RF-002), demostraciones sin colisiones (RF-003), análisis/pruebas EKF (RF-004), pruebas de activación fault stop (RF-005) y pruebas/inspección CSP (RF-006), usando criterios cuantitativos definidos. \cite{nasa_cm, segoldmine_ppi, omg_sysml_mbse}
\end{enumerate}

\subsection{Minimal Configuration and Change Control (Tailored)}\label{subsec:cm-tailored}

Esta subsección define un proceso mínimo de control de configuración y cambios (CM/CC) proporcional al contexto académico, para preservar consistencia entre: SyRS/SRS (Sección~11), modelo MBSE/SysML (BDD/IBD, comportamiento), PBS/ASM, matriz N2, ICDs y V\&V.

\paragraph{Configuration Items (CIs) mínimos.}
Los siguientes artefactos se consideran \textit{Configuration Items} (CIs) bajo control de cambios:
\begin{itemize}
  \item \textbf{CI-REQ}: SyRS/SRS (definición maestra de requisitos en Sección~11).
  \item \textbf{CI-MODEL}: Modelo SysML (BDD/IBD/Interfaces, Activity/State).
  \item \textbf{CI-PBS}: PBS/ASM + BOM COTS (mapeo función$\rightarrow$COTS).
  \item \textbf{CI-N2}: Matriz N2 de conectividad (interfaces físicas y lógicas).
  \item \textbf{CI-ICD}: ICD de comunicaciones (Rover--GCS) e ICD interno (buses/pines mínimos).
  \item \textbf{CI-VV}: Matrices de verificación y procedimientos VT-* (SyRS/SRS).
  \item \textbf{CI-CODE}: Repositorios de software (tags/releases) y configuración de campaña.
  \item \textbf{CI-DATA}: Datasets/ground truth y scripts de evaluación (si aplican a métricas de percepción).
\end{itemize}

\paragraph{Proceso mínimo de cambio (CR $\rightarrow$ Impacto $\rightarrow$ Aprobación $\rightarrow$ Baseline).}
Todo cambio propuesto DEBERÁ registrarse como una \textit{Change Request} (CR) con un identificador único (CR-xxx) y seguir el flujo:
\begin{enumerate}
  \item \textbf{Registro CR}: motivo, descripción, CIs afectados, prioridad, relación con UC-01 y riesgo.
  \item \textbf{Impact analysis}: identificar impacto en requisitos (CI-REQ), arquitectura/IBD/N2 (CI-MODEL/CI-N2), implementación física (CI-PBS), interfaces/ICD (CI-ICD) y V\&V (CI-VV).
  \item \textbf{Aprobación (mini-CCB)}: Responsable SE + líder I\&T + líder SW aprueban/rechazan. Si afecta requisitos, se actualiza CI-REQ (Sección~11) y su trazabilidad a V\&V.
  \item \textbf{Implementación}: aplicar cambios en los CIs correspondientes, referenciando CR-xxx en commits/notas.
  \item \textbf{Regresión V\&V}: ejecutar pruebas VT afectadas y actualizar evidencia (CI-VV).
  \item \textbf{Baseline update}: etiquetar la versión (tag/release), actualizar historial de revisiones y congelar baseline para campaña.
\end{enumerate}

\paragraph{Plantilla mínima de CR (recomendada).}
\begin{itemize}
  \item \textbf{CR-ID}: CR-xxx
  \item \textbf{Motivo}: (p.ej., inconsistencia, riesgo de integración, mejora V\&V)
  \item \textbf{CIs afectados}: CI-REQ/CI-MODEL/CI-PBS/CI-N2/CI-ICD/CI-VV/CI-CODE/CI-DATA
  \item \textbf{Impacto en requisitos}: IDs afectados (SyRS-xxx/SRS-xxx/RF/RNF/CON)
  \item \textbf{Impacto en V\&V}: VT-IDs afectados y cambio en criterios/procedimientos
  \item \textbf{Decisión mini-CCB}: Aprobado/Rechazado + fecha + responsables
\end{itemize}
